\subsection{Agent-native publishing and composite epistemic scores}

A recent architectural neighbor is Traxia, which proposes agent-native scientific publishing around verifiable epistemic artefacts carrying claims, reasoning traces, provenance, agent identity, replication records and contradiction handling \cite{dogah2026traxia}. This occupies substantial prior-art territory around machine-readable scientific artefacts, agent attribution, verifiable publishing and persistent contradiction records. Epistemic mechanics therefore does not claim that these functions arise uniquely from RAKL.

The comparison also exposes a substantive design fork. Traxia defines a composite Epistemic Confidence Score by weighting claim confidence, reproducibility and trace completeness. Such a scalar can be useful as a platform ranking policy. The theorem in Section~\ref{sec:authority} makes a narrower structural point: whenever epistemic states contain incomparable authority coordinates, no real-valued score can preserve the partial order by a biconditional order embedding. RAKL consequently treats any scalar ranking as a declared decision projection, not as the scientific authority state itself. This is not a claim that composite scores are invalid; it is a requirement that their policy semantics remain separate from the multidimensional evidence relation they summarize.

Traxia and RAKL also differ in where they place the release boundary. Traxia focuses on publication infrastructure, identity, review and reputation. Epistemic mechanics focuses on controlled scientific-state transition and on whether continuing discovery still changes the registered claim/evidence/novelty state. The two architectures are therefore adjacent and potentially composable: a bounded-saturated RAKL research state could be packaged into an agent-native publication artefact, while a publication record could provide provenance objects for later RAKL assimilation.
